\subsection{The exact dyadic point locus in the pointed Fermat-eight class}

\begin{statement}{A.189}{Proved}
{A rank-one logarithm calculation and a dyadic constraint on the full metabelian fibre}
Put
\[
 C:\ y^2=x^5-54x,\qquad J=\operatorname{Jac}(C),\qquad
 P=(3,9),\quad D=[P-\infty],\quad T=[(0,0)-\infty].
\tag{A.189.1}
\]
Use the actual completed coordinates of
(A.186.18)--(A.186.21):
\[
          \widehat{J(\Q)}_2=\Z_2D\oplus(\Z/2)T.
\]
The following local intersection
is exact:
\[
\left\{R\in C(\Q_2):\ x(R)\in3+8\Z_2,\quad
 \kappa_2([R-\infty])\in
 \operatorname{loc}_2\kappa_2(\widehat{J(\Q)}_2)\right\}
                   =\{(3,9),(3,-9)\}.
\tag{A.189.2}
\]
The corresponding global completed classes are respectively
\(D\) and \(-D\).

For the proper Fermat curve \(C_8:3u^8-2v^8=1\),
the exact morphism used below is
\[
 q:C_8\longrightarrow C,\qquad
       x=3u^2/v^2,\quad y=9u/v^5.
\tag{A.189.3}
\]
It gives
\[
                         C_8(\Q)=\{(\pm1,\pm1)\}.
\tag{A.189.4}
\]
Moreover every local point of \(C_8(\Q_2)\) whose
compact \(J\)-Kummer image is the localization of a
global class in \(\operatorname{Sel}(\Q,T_2J)\)
is one of these four points.

Let \(U=\mathbb P^1_\Q\setminus\{0,1,\infty\}\),
\(\Pi=\pi_1(U_{\overline\Q})\), and interpret both
derived subgroups as closed subgroups.  Every section
of the actual arithmetic extension
\(\pi_1(U)/\Pi''\to G_\Q\) that is locally
the image of a point or a cusp and has the finite
Kummer coordinates
\[
             (\kappa_8(t),\kappa_8(1-t))=([3],[-2])
\tag{A.189.5}
\]
has ordinary dyadic centre exactly \(t_2=3\).
After choosing its fixed sheet on
\(3u^8-2v^8=1\), its global compact two-adic image
in \(J\) is \(D\) or \(-D\), according to that
sheet.  In particular this conclusion applies to
every full Selmer section in the same finite class.
This specifies the dyadic centre and this compact
coefficient image; it asserts no equality of the full
global section with the point section at three.
\end{statement}

\begin{proof}
\emph{The abelian logarithm on the smooth local discs.}
Write \(f(x)=x^5-54x\), and use the basis
\(\omega_1=dx/y,\ \omega_2=x\,dx/y\).
Let \(\log:J(\Q_2)\to\Q_2^2\) be the abelian
logarithm evaluated on this basis.  We will also use
its scalar extension to one quadratic extension of
\(\Q_2\).  The logarithm is a homomorphism, and
on the kernel of reduction to the Neron special fibre
it is given by formal antiderivatives, including at
bad reduction; see
\cite[\S4.1, Remark 4.1, author version p.4]{LegMcCallumPoonenChabauty}.
Here is the precise disc comparison needed for our
equations.  On the affine integral model with \(x\)
a unit, \(f'(x)=5x^4-54\) is a unit.  This is a
smooth curve over the local valuation ring.  Its
Abel--Jacobi map extends to the Neron model by
\cite[Theorem 5.7.25 and Proposition 5.7.26, p.143]{LegPoonenRationalPoints}.
Two points with the same reduction therefore have
Abel--Jacobi difference in the formal group.  Pulling
back its invariant forms gives exactly
\(\omega_1,\omega_2\).  Thus a convergent power-series
primitive on their common smooth disc computes the
logarithm of that difference.

This comparison also follows directly on the formal
completions: invariant forms have integral coefficients
after one fixed scalar multiplication, and their
primitive coefficients have denominators bounded by
their total degrees.  Consequently the primitives
converge at every tuple of positive valuation.  The
formal chain rule identifies their pullbacks with the
curve primitives.  Dividing by the fixed scalar proves
the assertion for our chosen basis.  This argument
applies over the quadratic local extension as well.

\emph{A literal multiple of the global generator.}
We specify the rational divisor arithmetic in full.
For a monic quadratic \(U\) and linear \(V\) with
\(U\mid f-V^2\), write \([U,V]\) for the class
of its two points minus \(2\infty\), with
multiplicities.  The initial pair
\[
 U_2=(x-3)^2,\qquad V_2=(39x-99)/2
\tag{A.189.6}
\]
represents \(2D\): its linear polynomial has value
nine and derivative \(f'(3)/18\) at three.
The following rule gives the actual doubling relation.
In \(\Q[x]/(U)\), take the inverse of \(2V\), and put
\[
\begin{aligned}
 L&=V+U\left(\frac{f-V^2}{U}(2V)^{-1}\bmod U\right),\\
 U'&=\operatorname{monic}\left(\frac{f-L^2}{U^2}\right),
       &V'&=-L\bmod U'.
\end{aligned}
\tag{A.189.7}
\]
For the two inputs used here the inverse exists and
\(L\) has degree three, while \(U'\) has degree two.
The identity \(f-L^2=\lambda U^2U'\) for a
nonzero rational \(\lambda\) gives the divisor
of \(y-L\): twice the two old points plus the
two new points, minus \(6\infty\).  Negating the
new points gives \([U',V']=2[U,V]\).
Thus the formula proves the group comparison rather
than presupposing an arithmetic routine.

Applying (A.189.7) twice gives
\[
\begin{aligned}
 U_4&=x^2-\frac{5132754}{769129}x
                         +\frac{10169721}{769129},\\
 V_4&=\frac{44708953611x-121694660451}{2698104532},\\
 U_8&=x^2+a_8x+b_8,\\
 a_8&=\frac{21417625950097970853055914846}{16459917593113792410798392089},\\
 b_8&=\frac{308227194535288413050444696841}{16459917593113792410798392089}.
\end{aligned}
\tag{A.189.8}
\]
Define \(V_8\) by the second application of
(A.189.7), so all of its rational coefficients are
specified without a choice.  Set
\[
 c=-\tfrac12[x]U_8,\qquad d=c^2-U_8(0),\qquad
 V_8(c+t)=v_0+v_1t,\qquad F=f(c).
\tag{A.189.9}
\]
Direct rational arithmetic gives the exact data
\[
\begin{gathered}
 c\equiv249\pmod{256},\qquad
 (v_2(c),v_2(d),v_2(F),v_2(v_0),v_2(v_1))
                         =(0,6,0,2,-3),\\
 v_2\!\left(\frac{F}{v_1^2d}-1\right)=4.
\end{gathered}
\tag{A.189.10}
\]
Both polynomial divisor identities and every entry of
this valuation display are checked as exact rational
identities in the retained certificate.

Choose \(\delta^2=d\), and let
\(P_\pm=(c\pm\delta,v_0\pm v_1\delta)\).
Then
\(8D=[P_++P_--2\infty]=[P_+-\iota P_-]\),
where \(\iota(x,y)=(x,-y)\).
There is a unique \(k\in1+8\Z_2\) with
\(k^2=F/(v_1^2d)\).  Put
\(\beta=v_1\delta k\), so \(\beta^2=F\).
The branch
\[
                  y(t)=\beta\left(\frac{f(c+t)}{F}\right)^{1/2}
\tag{A.189.11}
\]
converges for \(v_2(t)\geq3\).  At \(t=\delta\)
it equals \(v_0+v_1\delta\), and at \(t=-\delta\)
it equals \(v_1\delta-v_0\).  Indeed all four
expressions are roots of the appropriate curve
equations, and each is congruent to \(v_1\delta\)
modulo four; its negative has difference of valuation
one.  This proves the branch identities.  The two
endpoints have the same smooth reduction, since their
coordinate differences have positive valuation.
The formal comparison above therefore computes
\(\log(8D)\) by integrating (A.189.11) from
\(-\delta\) to \(\delta\).

\emph{The slope of the global logarithm.}
Write
\[
 \left(\frac{f(c+t)}{F}\right)^{-1/2}
                        =\sum_{n\geq0}b_nt^n.
\tag{A.189.12}
\]
Here \(b_0=1\), and for \(n\geq1\),
\(v_2(b_n)\geq1-2n\).
To prove the bound, the nonconstant polynomial inside
the binomial expansion has integral coefficients and
zero constant term.  In degree \(n\), only the
binomial terms with index at most \(n\) occur.
For every positive index \(j\),
\(\binom{-1/2}{j}=(-1)^j\binom{2j}{j}/4^j\)
has valuation at least \(1-2j\), since the central
binomial coefficient is even.  This proves the stated
bound for every coefficient and proves convergence
of the series and integrals used here.

Define convergent rational two-adic series
\[
 S=\sum_{j\geq0}\frac{b_{2j}d^j}{2j+1},\qquad
 N=\sum_{j\geq0}\frac{b_{2j+1}d^{j+1}}{2j+3}.
\tag{A.189.13}
\]
The terms of \(S-1\) have valuations at least
\(2j\) for \(j\geq1\), and those of \(N\)
have valuations at least \(2j+4\) for \(j\geq0\).
Their denominators are odd.  Hence
\(S\in1+4\Z_2\), \(N\in16\Z_2\), and
the exact logarithm is
\[
 \log(8D)=\frac{2\delta}{\beta}(S,cS+N),\qquad
       \rho:=\frac{\log(D)_2}{\log(D)_1}
                      =c+N/S\equiv1\pmod4.
\tag{A.189.14}
\]
In particular its first coordinate is nonzero; it
has valuation four for \(8D\) and valuation one
for \(D\).  The expression is rational over
\(\Q_2\), since \(\delta/\beta=1/(v_1k)\).

\emph{All points in the relevant curve disc.}
Let \(R=(3+h,y)\), where \(h\in8\Z_2\).
Choose \(y_0\in\{9,-9\}\) with
\(y\equiv y_0\pmod4\), and put
\(P_0'=(3,y_0)\).  The unique branch through
\(P_0'\) is
\(y(t)=y_0(f(3+t)/81)^{1/2}\) for
\(t\in8\Z_2\).  The endpoint \(y(h)\)
equals the chosen \(y\), by the same root comparison
modulo four.  Write
\((f(3+t)/81)^{-1/2}=\sum a_nt^n\).
The preceding binomial proof gives
\(a_0=1\) and \(v_2(a_n)\geq1-2n\) for
every \(n\geq1\).

For \(h\ne0\), let \((I_1,I_2)\) be the
logarithm of \([R-P_0']\).  The same smooth-disc
comparison gives
\[
\begin{aligned}
 I_1&=\frac{h}{y_0}
       \sum_{n\geq0}\frac{a_nh^n}{n+1},\\
 I_2-3I_1&=\frac{h^2}{y_0}
       \sum_{n\geq0}\frac{a_nh^n}{n+2}.
\end{aligned}
\tag{A.189.15}
\]
Writing \(e=v_2(h)\geq3\), each term after the
constant in the first sum has valuation at least
\((e-2)n+1-v_2(n+1)\geq1\).
That sum is therefore in \(1+2\Z_2\), and
\(I_1\ne0\).  In the second sum the constant
is \(1/2\), while every later term has valuation
at least \((e-2)n+1-v_2(n+2)\geq1\).
Consequently
\[
                    \frac{I_2}{I_1}\equiv3\pmod4.
\tag{A.189.16}
\]
The logarithm of the completed global Mordell--Weil
image lies in \(\Q_2\log D\): the torsion class
has logarithm zero and (A.186.21) is the exact
completed generator map.  The local logarithm
factors continuously through the local two-primary
completion.  For example this follows by taking a
small open formal subgroup on which logarithm is
an isomorphism onto a \(\Z_2\)-lattice, and then
using its finite quotient.  Thus equality of the
local continuous Kummer classes implies equality
of these logarithms.  Subtracting the rational
class \([P_0'-\infty]=\pm D\), (A.189.14) and
(A.189.16) contradict membership in that line
unless \(h=0\).  This proves (A.189.2).

It also proves the asserted uniqueness of its
global completed classes.  Write such a class as
\(aD+eT\).  The nonzero logarithm forces
\(a=1\), respectively \(-1\).  The local
two-adic Kummer class of \(T\) is nonzero: otherwise
\(T\) would have compatible divisibility by
arbitrarily large powers of two in \(J(\Q_2)\).
Each such preimage would be two-primary torsion
of unbounded order, whereas an open logarithm
subgroup is torsion free and its quotient is finite.
Hence \(e=0\).

\emph{The actual Fermat map and its dyadic points.}
Substitution in (A.189.3) gives
\(y^2=x^5-54x\) exactly; this is the first
genus-two map of A.182 followed by the monic
isomorphism of A.186.  The rational map extends
to the proper curves.
There is no dyadic point at infinity on \(C_8\),
since it would make \(2/3\) an eighth power.
Neither \(u=0\) nor \(v=0\) gives a dyadic
point: the first has a valuation not divisible
by eight, and the second would make \(1/3\)
an eighth power, contrary to its unit residue
modulo thirty-two.

For any remaining dyadic point, put
\(a=v_2(u),b=v_2(v)\).  The minimum of
\(8a,1+8b,0\) must occur at least twice.
The first two are incongruent modulo eight,
so necessarily \(a=0\) and \(b\geq0\).
Since an odd eighth power is one modulo
thirty-two, the equation now forces \(b=0\)
as well.  Thus \(u,v\) are units and
\(x=3(u/v)^2\equiv3\pmod8\).
Applying (A.189.2) gives \((u/v)^2=1\).
The Fermat equation then gives \(v^8=1\),
whose rational dyadic roots are exactly \(\pm1\).
Hence \(u,v\in\{1,-1\}\).
This proves the local assertion and (A.189.4),
because a rational point supplies an actual
global Mordell--Weil class.  For a global
two-adic Selmer class the same argument applies
by (A.186.2) and the finite-to-continuous
comparison of A.187.

\emph{The full metabelian section comparison.}
The affine cover
\[
 V:\ 3u^8-2v^8=1,\quad uv\ne0,
 \qquad t=3u^8,\quad1-t=-2v^8
\tag{A.189.17}
\]
has geometric subgroup
\(H=\ker(\Pi\to(\Z/8)^2)\).
It is exactly the twist of the two eighth-root
torsors with the finite coordinates in
(A.189.5), and \((1,1)\) is its rational
sheet above three.  Conjugating by a lift of
an abelian finite conjugator makes the finite
arithmetic section literally the one defined
by that sheet.  Since \(\Pi'\subset H\),
we have \(\Pi''\subset\overline{[H,H]}\).
The given metabelian section consequently
maps into the actual cover extension
\[
 \pi_1(V)/\overline{[H,H]}\longrightarrow G_\Q.
\tag{A.189.18}
\]
Compactification, (A.189.3), and the
Abel--Jacobi map with origin \(\infty\)
give the equivariant geometric homomorphism
\(H^{\rm ab}\to T_2J\), together with
its arithmetic extension map.  At every
finite multiplication level the resulting
section class is a global Kummer Selmer
class.  To see the local condition, an
ordinary local image and its fixed sheet
give the actual point in the fibre of
(A.189.17).  A cuspidal image gives its
rational boundary point by the finite-sheet
decomposition-group comparison of A.178;
compactification kills its horizontal
inertia.  The same fixed-sheet argument
at a real place gives the corresponding
real component.  Thus these maps have
exactly the local Kummer images used in
(A.186.2), at all places and all levels.

At two the proper boundary of (A.189.17)
has no rational dyadic point, as just
proved.  The local lift is therefore an
ordinary point of \(V(\Q_2)\).
The preceding local calculation makes it
one of \((\pm1,\pm1)\); its image on
\(U\) is exactly three.  The global
compact coefficient class is \(D\) or
\(-D\) by the already proved uniqueness.
This proves every stated section assertion.

The exact divisor and series arithmetic is
retained in
\path{code/point3_genus2_chabauty_certificate.py},
with its output and rational data.  It checks
both doubling identities, every valuation in
(A.189.10), fourteen symmetric-series
coefficients, and 128 ordinary-disc tests.
The binomial estimates above control all
remaining coefficients; no finite truncation
is being used as a proof of an infinite
series identity.
\end{proof}
